\section{Síntesis y ampliación}

Las dos conferencias explican la Transformer-like computation desde direcciones distintas. La ruta de SDM parte de la geometría de alta dimensión: el intersection count de las distributed addresses aproxima una similitud exponencial, de donde se obtiene la softmax retrieval; el cerebellar circuit ofrece una implementación candidata de read/write. La ruta de TEM parte de las necesidades cognitivas: separa la reusable structure del changing content, los vincula mediante una fast associative memory y luego reformula la memory read como modern Hopfield attention.

\subsection{Comparación unificada de las dos rutas}

La siguiente tabla no fusiona el cerebellum y el hippocampus en un solo modelo cerebral; compara qué parte de la attention explica cada uno. La primera conferencia se centra más en el weighting kernel y en la geometría del almacenamiento; la segunda, en el positional/structural state, el rapid binding y la compositional inference. Ambas enfatizan la key/content separation, la distributed memory y la uncertainty about biological implementation.

\begin{longtable}{@{}L{0.18\textwidth}L{0.24\textwidth}L{0.25\textwidth}L{0.21\textwidth}@{}}
\toprule
Ruta & Pregunta principal & Correspondencia en el Transformer & Evidencia más sólida \\
\midrule
SDM geometry & Cómo una noisy query recupera recuerdos según su similitud & softmax attention weights & Aproximación de intersecciones en alta dimensión y experimentos numéricos \\
Cerebellar mapping & Cómo un circuito neuronal separa activación, escritura y lectura & keys/values y pooled retrieval & Factibilidad anatómica y experimentos candidatos \\
Cognitive map / TEM & Cómo reutilizar la estructura relacional y vincular rápidamente contenido nuevo & recurrent position state + memory & Comparación entre conducta, modelos y representaciones neuronales \\
Modern Hopfield TEM & Cómo ampliar la capacidad de la associative read e integrarla con el context & query/key/value attention & Correspondence de fórmulas y beneficios de implementación \\
\bottomrule
\end{longtable}

\begin{importantbox}{El verdadero método «neuroscience-inspired» del curso}
No se trata de copiar el nombre de alguna región cerebral, sino de identificar primero el problema computacional, proponer después un modelo con un inductive bias explícito y, por último, establecer vínculos falsables entre las matemáticas, la conducta en la tarea, las representaciones internas, los datos neuronales y los experimentos de intervención. Una brain analogy vale más que una metáfora solo cuando genera predicciones nuevas.
\end{importantbox}

\subsection{Seis lecciones para la investigación sobre el Transformer}

\begin{enumerate}
  \item \textbf{Explicar las condiciones geométricas del softmax.} Revisar la norm, la temperature, la dimension y la approximation region, en lugar de tratar el softmax como un valor predeterminado misterioso.
  \item \textbf{Distinguir la pattern memory de corto plazo de la distributed memory de largo plazo.} La context precision y la parameter capacity enfrentan interferencias y grados de trazabilidad distintos.
  \item \textbf{Dejar explícita la separación entre key y value.} La señal de coincidencia y el contenido devuelto son distintos; este es el mecanismo común a la attention, la SDM y el hippocampal binding.
  \item \textbf{Hacer que la representación de la posición refleje la estructura de la tarea.} El sequence index es solo la estructura más simple; la graph transition, la action y el relational state pueden aportar positional dynamics más potentes.
  \item \textbf{Reservar un fast state para el vínculo rápido.} El contenido de un entorno nuevo no siempre debe depender de actualizaciones lentas de parámetros; la memory/fast weights puede encargarse del episode-specific knowledge.
  \item \textbf{Convertir la biological plausibility en experimentos.} Especificar qué cells registrar, cuándo se escribe, cómo intervenir y qué resultado refutaría el mapping.
\end{enumerate}

\begin{warningbox}{Las tres afirmaciones de esta clase que más fácilmente se exageran}
«Attention approximates SDM» no significa que la attention y la SDM sean idénticas con cualquier parámetro; «el cerebelo puede implementar la SDM» no significa que esté demostrado que el cerebelo ejecuta un Transformer; «TEM equivale aproximadamente a un Transformer» no significa que el hippocampus sea un modelo de lenguaje. Cada afirmación debe acompañarse del nivel de comparación, los supuestos y el espacio de contraejemplos.
\end{warningbox}

\subsection{Lecturas adicionales}

\begin{itemize}
  \item Bricken and Pehlevan, \href{https://proceedings.neurips.cc/paper/2021/hash/8171ac2c5544a5cb54ac0f38bf477af4-Abstract.html}{\emph{Attention Approximates Sparse Distributed Memory} (NeurIPS)}
  \item Whittington et al., \href{https://www.sciencedirect.com/science/article/pii/S009286742031388X}{\emph{The Tolman--Eichenbaum Machine} (Cell)}
  \item Whittington et al., \href{https://arxiv.org/abs/2112.04035}{\emph{Relating Transformers to Models and Neural Representations of the Hippocampal Formation} (arXiv)}
  \item Bricken et al., \href{https://openreview.net/forum?id=JknGeelZJpHP}{\emph{Sparse Distributed Memory is a Continual Learner} (OpenReview)}
  \item Behrens et al., \href{https://www.nature.com/articles/s41593-022-01153-y}{\emph{How to Build a Cognitive Map} (Nature Neuroscience)}
\end{itemize}

Al leer, conviene aplicar siempre la misma lista de verificación: cuál es el objeto computacional, cómo se factoriza la representación, cómo se escribe y se lee la memory, qué equivalencias son exact/approximate, qué resultados provienen del modelo y cuáles de los neural data, y qué experimento podría refutar el mapping. Solo así la neuroscience se convierte en diseño de modelos y en predicciones científicas, en lugar de limitarse a ponerle etiquetas de regiones cerebrales a arquitecturas existentes.

\end{document}
